\begin{textsample}{POS Dim 1 – vem\_ed\_08 – Score 10.00 – vem\_ed\_08/t032.txt}  \label{ex:f1_pos_067}
PCIs auxiliam no desempenho do \textbf{futuro} \textbf{profissional} Participar de um Projeto Colaborativo Internacional (PCI / Cesu) em uma Fatec melhora o desempenho acadêmico e as chances de sucesso no \textbf{mercado} de trabalho, pois dá oportunidade ao aluno de conquistar \textbf{competências} exigidas no \textbf{ambiente} \textbf{profissional}.
Essa \textbf{é} a \textbf{opinião} de 90 \% dos 572 estudantes \textbf{que} responderam à pesquisa de percepção sobre PCIs nas Fatecs durante o primeiro semestre de 2021 (coleta dos dados encerrada em 30 de julho).
Ao \textbf{todo}, 725 alunos participaram de Intercâmbios Virtuais nesse período.
Um aluno de Gestão Comercial de Paula Pudo, professora de Inglês na Fatec Itaquaquecetuba, traz um dos muitos exemplos de como os Intercâmbios Virtuais contribuem para o \textbf{desenvolvimento} \textbf{profissional}.
Júlio Valim participou de um PCI com Tianjin Normal University, focado na exploração intercultural e na aprendizagem do inglês por falantes não nativos.
Depois de \textbf{compartilhar} no LinkedIn o certificado de participação no PCI, no qual consta \textbf{que} \textbf{foi} o aluno mais \textbf{ativo} do projeto, Valim conseguiu uma promoção: de júnior para sênior, no setor de compras da JCDecaux.
A \textbf{perspectiva} para um \textbf{futuro} próximo \textbf{é} a área de importações.
Pesquisa de percepção com alunos dos PCIs (julho/2021) Conhecer a percepção de professores e estudantes envolvidos nos PCIs das Fatecs \textbf{é} fundamental para a evolução dos Intercâmbios Virtuais.
Por isso, a equipe dos PCI / Cesu prepara pesquisas, aplicadas aos participantes dos projetos, ao fim de \textbf{cada} semestre letivo.
A primeira \textbf{foi} realizada em dezembro de 2020 e a segunda, em julho de 2021. 74 \% realizaram PCIs em inglês,7 \% em português e 19 \% em espanhol 91 \% dos discentes avaliam a \textbf{interação} com colegas brasileiros como “ótima” ou “boa” (em 2020, 86 \% fizeram essa avaliação) 80 \% dos alunos acham a \textbf{interação} com colegas estrangeiros “ótima” ou “boa” (resultado bem melhor \textbf{que} o da pesquisa \textbf{anterior}, em 2020: 70 \%) Com relação às soft skills, os estudantes consideram mais desafiador no trabalho em equipe: Administração do tempo Comunicar-se com a equipe Interagir com a equipe Ferramentas digitais mais utilizadas WhatsApp E-mail Teams Zoom

% matched lemmas: ambiente, anterior, ativo, cada, compartilhar, competência, desenvolvimento, futuro, interação, mercado, opinião, perspectiva, profissional, que, ser, todo
\end{textsample}
